\section{Limitazioni e discussione}
\label{sec:limiti}
\subsection{Persistenza e guasti}
La persistenza del log Kafka non implica il ripristino dello stato delle
finestre. Worker e Global mantengono gli aggregati in RAM e committano gli
offset dopo ogni elaborazione; non esiste un checkpoint coordinato fra
offset e stato. Un crash può quindi perdere contributi già consumati e
committati, che non verrebbero riletti al riavvio. Non è garantita una
semantica exactly-once end-to-end. Il sink PostgreSQL finale non è uno
state store degli aggregatori.

Per questo motivo un rebalance dopo l'inizio dell'elaborazione invalida la
run; il worker rifiuta inoltre la ripartenza da offset già committati.
Le repliche devono essere pronte prima del replay e restare stabili fino al
termine. Un errore richiede una nuova esecuzione completa con topic e offset
puliti. Il broker Kafka è singolo, con fattore di replica uno: gli
acknowledgment richiesti non realizzano alta disponibilità. La tolleranza ai
guasti non è il requisito individuale selezionato.

\subsection{Completezza, carico e rappresentatività}
QoS 0 e code limitate possono perdere telemetria. Il protocollo EOS certifica
il completamento del flusso accettato dai componenti, non il recupero degli
eventi persi a monte. Le run senza perdite devono quindi essere identificate
tramite riconciliazione dei contatori e controlli sui risultati. Negli
esperimenti di saturazione o rete degradata le perdite sono a loro volta
un risultato da riportare; un throughput apparentemente stabile ottenuto
scartando dati non soddisfa il requisito di scalabilità.

L'attesa di tutte le sorgenti evita risultati prematuri ma può trattenere
finestre in memoria quando un Edge non avanza. Un timeout del runner
interrompe una prova incompleta senza autorizzare il Global a ignorare la
sorgente. Non sono implementati checkpoint, migrazione dello stato o una
politica di esclusione dinamica degli Edge.

% Bilanciamento manuale dell'ultima pagina della bozza; rimuovere se cambia il testo.
\newpage
Il workload bilancia i sensori per zona e riguarda un solo tipo di dispositivo
in un solo mese. Le medie sono pesate per numero di campioni validi: sensori
più frequenti contribuiscono maggiormente. I risultati non possono essere
generalizzati automaticamente ad altri dataset, algoritmi più costosi o
installazioni geograficamente distribuite. Inoltre, la condivisione degli
host e i limiti per container introducono contesa; la scalabilità delle
repliche entro un host non dimostra un'illimitata capacità di espansione.

\section{Conclusioni}
L'applicazione realizza una pipeline Edge--Cloud in Go con validazione locale,
aggregazione gerarchica e comunicazione persistente nei passaggi Kafka. La
separazione fra partizioni logiche ed esecutori mantiene invariato il
significato degli output al variare delle repliche. Progresso ed EOS
permettono di concludere le finestre soltanto quando ogni contributo è
determinato, includendo le sorgenti senza dati.

Le verifiche funzionali e gli artefatti preliminari forniscono una base per
la valutazione, ma non consentono ancora di concludere che la scalabilità
richiesta sia dimostrata. Il completamento del lavoro richiede una campagna
EC2 valida e ripetuta, l'analisi dei tempi di disponibilità dei risultati e
la valutazione di latenza e banda emulate. Le conclusioni quantitative
dovranno essere aggiornate dopo tali esperimenti, conservando espliciti i
limiti dello stato volatile e della telemetria con possibili perdite.
